\chapter{Objetivos y alcance}
\label{chap:objetivos}

Este capítulo delimita el propósito y la extensión del trabajo. En primer
lugar, se formula el objetivo general y se concreta en una serie de objetivos
específicos. A continuación, se describe el alcance de la solución y se señalan
los aspectos que quedan fuera de él. Finalmente, se establecen los criterios
que se utilizarán para evaluar la viabilidad y la reutilización del enfoque
propuesto.

\section{Objetivo general}
\label{sec:objetivo-general}

El objetivo general de este Trabajo Fin de Máster es diseñar e implementar un
entorno para crear lenguajes de dominio específico con sintaxis basada en
bloques. El enfoque propuesto consiste en describir el dominio a alto nivel
mediante un metamodelo Ecore anotado o mediante una descripción textual
\texttt{.m2b}, transformar esa entrada en un modelo intermedio EMF y generar a
partir de esa representación un editor Blockly ejecutable. La ruta textual
\texttt{.m2b} se entiende como un DSL de definición de lenguajes: la usa el
diseñador del lenguaje para especificar el editor, mientras que el usuario final
interactúa con el editor Blockly generado.

Este objetivo combina dos enfoques. Por un lado, Blockly ofrece la infraestructura para
crear editores visuales, pero el desarrollador debe definir los bloques, campos,
conexiones y generadores \cite{blocklyWhatIs}. Por otro lado, la ingeniería
dirigida por modelos usa modelos como artefactos centrales para especificar y
generar software \cite{brambilla2017modelDriven}. El TFM conecta ambas ideas para
reducir el trabajo manual necesario en cada nuevo lenguaje.

\section{Objetivos específicos}
\label{sec:objetivos-especificos}

Para alcanzar el objetivo general se han definido los siguientes objetivos
específicos:

\noindent\textbf{OE1. Analizar la estructura de los lenguajes basados en
bloques.} Identificar qué necesita un editor de bloques para representar un
dominio: tipos de bloques, campos, categorías, conexiones, restricciones,
serialización y, cuando sea necesario, generación de código. Este objetivo se
justifica por la variedad de diseños que existen en la programación por bloques
\cite{lin2021landscape}.

\medskip
\noindent\textbf{OE2. Definir una ruta de generación basada en metamodelos
Ecore.} Permitir que un metamodelo de dominio actúe como fuente principal de
generación. Las clases se mapean a tipos de bloque, los atributos a campos, las
referencias de contención a entradas estructurales y las referencias no
contenedoras a mecanismos de selección o enlace entre elementos.
Las cardinalidades se traducen en comprobaciones sobre cuántos elementos son
obligatorios o están permitidos. Esta ruta representa el enfoque
metamodelo-céntrico del trabajo: generar código Blockly a partir de un
metamodelo, posiblemente anotado con información adicional sobre su sintaxis
visual basada en bloques.

\medskip
\noindent\textbf{OE3. Proporcionar una ruta textual \texttt{.m2b} para definir
lenguajes de bloques.} Como complemento a la ruta Ecore, se define un DSL
textual implementado con Xtext para describir dominios de forma compacta. Su
función es ofrecer una notación de autoría para diseñadores de lenguajes, capaz
de declarar categorías, clases, atributos, contenciones, referencias,
validaciones y metadatos de interfaz, y normalizar después esa información hacia
el mismo modelo intermedio. Xtext es adecuado para esta tarea porque genera
infraestructura de parsing, modelo abstracto basado en EMF y soporte de edición
a partir de una gramática
\cite{xtextProject,bettini2013xtext}.

\medskip
\noindent\textbf{OE4. Diseñar un modelo intermedio independiente de la salida.}
Separar la transformación del dominio de la generación de Blockly. Para ello, la
ruta Ecore y la ruta textual \texttt{.m2b} se normalizan en una representación
común \texttt{EditorSpec}. Así, el generador no depende de los detalles de
lectura de Ecore ni de la sintaxis textual y se pueden añadir nuevas salidas con
menos cambios.

\medskip
\noindent\textbf{OE5. Generar editores Blockly ejecutables.} Implementar un
generador que produzca los ficheros necesarios para abrir un editor en el
navegador: definición de bloques, paleta, generadores, validaciones y página
HTML. El editor generado permite crear programas o modelos usando la
sintaxis visual de bloques. El estado del workspace puede guardarse en JSON y
cargarse posteriormente en el editor. También es posible exportar e importar la
representación JSON del modelo. Además, el contenido puede exportarse como XMI
de instancia del dominio para su procesamiento con herramientas compatibles con
EMF, aunque actualmente este XMI no puede cargarse de nuevo en el editor. La
generación también debe producir artefactos revisables, como el XMI intermedio,
un modelo de ejemplo y un informe de generación.

\medskip
\noindent\textbf{OE6. Incorporar mecanismos genéricos de validación, referencias
y personalización.} Un editor generado resulta más útil si no se limita a
mostrar bloques aislados. Por ello, el entorno incorpora comprobaciones para
detectar modelos incompletos, referencias entre elementos del modelo, opciones
de presentación y anotaciones que permiten adaptar la apariencia del editor al
dominio. Además, las reglas de validación se materializan como una vista Blockly
separada, de modo que puedan inspeccionarse visualmente y, en los casos
soportados, aplicarse de vuelta a la fuente.

\medskip
\noindent\textbf{OE7. Evaluar la viabilidad mediante un ejemplo conductor y
pruebas.} La solución se evalúa con AppMaker descrito por las dos rutas: como
metamodelo Ecore con anotaciones y como fichero textual \texttt{.m2b}. Este caso
permite analizar si ambas entradas pueden converger en el mismo modelo
intermedio y generar un editor Blockly operativo con artefactos revisables. Las
pruebas automáticas cubren el modelo intermedio, los adaptadores y los
generadores.

\section{Alcance de la solución}
\label{sec:alcance-funcional}

El alcance de la solución se organiza alrededor de una cadena generativa:
describir el dominio, transformarlo en una representación intermedia y generar
un editor Blockly ejecutable. La entrada puede ser un metamodelo Ecore anotado o
un fichero textual \texttt{.m2b}. En ambos casos, el sistema construye una
especificación intermedia y genera una implementación
HTML/JavaScript que puede abrirse en un navegador.

La Tabla \ref{tab:alcance-funcional} resume las dimensiones consideradas en la
solución. Estas dimensiones no se organizan como una lista aislada de funciones,
sino como los elementos necesarios para sostener el argumento central del
trabajo: que un dominio descrito a alto nivel puede transformarse de forma
sistemática en una sintaxis visual basada en bloques.

\begin{table}[h]
  \centering
  \begin{tabular}{|p{0.32\textwidth}|p{0.56\textwidth}|}
    \hline
    Dimensión & Papel en la solución \\
    \hline
    Entrada Ecore & Uso de clases, atributos, referencias, herencia,
    cardinalidades y anotaciones como información de partida para construir el
    editor visual. \\
    Ruta textual \texttt{.m2b} & Sintaxis Xtext para describir dominios, categorías,
    clases, atributos, referencias, contenciones, validaciones y plantillas de
    código con una notación más compacta, normalizable al mismo modelo
    intermedio. \\
    Modelo intermedio & Representación común de bloques, campos, categorías,
    referencias, entradas de valor, validaciones y metadatos de interfaz,
    independiente de la ruta de entrada. \\
	    Generación Blockly & Producción de bloques, toolbox, generadores,
	    validaciones y páginas HTML ejecutables a partir de la representación
	    intermedia. \\
	    Integración Eclipse & Plugins Eclipse/Xtext para editar ficheros
	    \texttt{.m2b}, ejecutar la generación desde el menú contextual y
	    publicar la herramienta como feature instalable desde un update site p2. \\
	    Artefactos trazables & Generación del XMI intermedio, informe de generación,
	    modelo de ejemplo, workspace visual de validaciones y documentación breve de
	    la salida producida. \\
	    Validaciones & Derivación de comprobaciones sobre campos obligatorios,
	    mínimos y máximos de elementos, unicidad, referencias obligatorias,
	    expresiones simples y restricciones acotadas de orden. \\
    Referencias & Representación de relaciones de contención y no contención,
    incluyendo listas desplegables dinámicas para seleccionar elementos del
    modelo. \\
    Personalización & Uso de colores, etiquetas, categorías, tooltips, widgets
    y metadatos de interfaz para adaptar la sintaxis visual al dominio. \\
    Evaluación & Uso de AppMaker por la ruta Ecore y por la ruta textual
    \texttt{.m2b}, junto con pruebas automáticas para
    aportar evidencia sobre adaptadores, generadores y modelo intermedio. \\
    \hline
  \end{tabular}
  \caption{Dimensiones incluidas en el alcance de la solución}
  \label{tab:alcance-funcional}
\end{table}

\section{Límites del trabajo}
\label{sec:limites}

Los límites del trabajo se derivan de su orientación generativa. El objetivo no
es construir una plataforma comercial completa ni un entorno final para un único
dominio. El objetivo es mostrar una cadena reutilizable para crear lenguajes de
dominio específico con sintaxis basada en bloques.

Por lo tanto, quedan fuera del alcance principal:

\begin{itemize}
  \item la importación genérica de modelos XMI existentes a un workspace
  Blockly;
  \item la integración completa de lenguajes de restricciones complejos como OCL;
  \item la sincronización automática de referencias bidireccionales complejas;
  \item la creación de interfaces de usuario finales altamente especializadas
  para cada dominio;
  \item la generación de código con análisis semántico avanzado o formateadores
  específicos de lenguajes de propósito general.
\end{itemize}

Algunos de estos elementos se exploran parcialmente, como la exportación a JSON
o XMI, la sincronización de referencias opuestas no contenidas, la traducción de
un subconjunto sencillo de invariantes OCL, la edición visual de reglas de
validación y la generación de código mediante plantillas. Se consideran
extensiones del objetivo principal y no equivalen a un entorno completo de
modelado, validación y ejecución.
Por lo tanto, se retoman explícitamente como líneas de trabajo futuro en el
Capítulo \ref{chap:conclusiones}, donde se discuten como ampliaciones posibles
de la herramienta y no como requisitos de la evaluación actual.

\section{Criterios de evaluación}
\label{sec:criterios-evaluacion}

La evaluación del trabajo se apoya en los siguientes criterios, orientados a
valorar si el enfoque propuesto es técnicamente viable y reutilizable:

\begin{enumerate}
  \item permitir definir distintos dominios mediante metamodelos Ecore con
  estructuras y características diferentes;
  \item generar automáticamente editores Blockly funcionales a partir de dichos
  metamodelos, sin introducir modificaciones específicas en el generador para
  cada dominio;
  \item permitir complementar la generación con información de presentación,
  como etiquetas, colores o categorías;
  \item ofrecer una sintaxis textual con Xtext para definir lenguajes de bloques
  sin editar directamente Ecore en casos sencillos o ejemplos;
	  \item generar código HTML/JavaScript que pueda ejecutarse en un navegador;
	  \item producir artefactos de trazabilidad, como XMI intermedio, informe de
	  generación y ficheros de validación revisables;
	  \item demostrar la reutilización de la arquitectura con un dominio
	  representativo descrito por la ruta Ecore y por la ruta textual
	  \texttt{.m2b}, normalizadas al mismo modelo intermedio;
	  \item disponer de pruebas, verificaciones y un mecanismo de instalación que
	  reduzcan el riesgo de que la solución funcione únicamente para un caso
	  preparado manualmente.
\end{enumerate}

Estos criterios responden a la motivación inicial del trabajo: evitar que cada
lenguaje basado en bloques requiera escribir de nuevo la infraestructura Blockly
de bajo nivel. Si el dominio puede describirse como metamodelo Ecore o como
modelo textual \texttt{.m2b}, el sistema lo transforma a \texttt{EditorSpec} y
genera un editor operativo, entonces el enfoque aporta evidencia frente a la
programación manual de cada editor.
